%!TEX root = ../report.tex
\documentclass[report.tex]{subfiles}
\begin{document}

    \chapter{Additional Implementation Details}
    
    This chapter provides additional implementation details of the \gls{cvo} approach. It includes the derivation of the measurement Jacobian, the numerical Jacobian calculation, and a comparison of contact estimation methods. 
    
    \section{Measurement Jacobian Derivation} \label{appendix:MeasurementJacobianDerivation}
    
    This section provides the derivation of the measurement Jacobian presented in Section~\ref{sec:measurementModelAndJacobian}. The Jacobian is derived by calculating the partial derivatives of the measurement function $h(x)$ with respect to each error-state variable. This analytical derivation used to get the Jacobian matrix is derived with AI assistance using Gemini CLI~\cite{geminicli2026}. The presented equations in the following sections are generated using the AI assistant ChatGPT~\cite{openai_chatgpt_jacobian_2026}, and they have been manually verified for correctness. 
    
    \subsection*{Rotation Matrix}

    The rotation matrix block of the measurement Jacobian is derived from the measurement function $h(x)$, which is defined as follows:

    \begin{equation}
    h(x) =
    -\left(
    [R_{body}^{odom}]^{T}v_{body}^{odom}
    +
    \left(\omega_{body}^{body}\times r_{pc}^{body}\right)
    \right)
    \end{equation}

    For a small rotation error $\delta R$, the first-order approximation of the body velocity is
        
    \begin{equation}
        [R_{body}^{odom}]^{T}v_{body}^{odom}
        \approx
        [R_{body}^{odom}]^{T}v_{body}^{odom}
        +
        \left[
        [R_{body}^{odom}]^{T}v_{body}^{odom}
        \right]_{\times}\delta R
    \end{equation}

    Substituting this approximation into the measurement function $h(x)$, we get,

    \begin{align}
        h(x,\delta R)
        &\approx
        -\left(
        [R_{body}^{odom}]^{T}v_{body}^{odom} + \left[ [R_{body}^{odom}]^{T}v_{body}^{odom} \right]_{\times}\delta R +
        \left(
        \omega_{body}^{body}\times r_{pc}^{body}
        \right)
        \right)
        \\
        &=
        -\left(
        [R_{body}^{odom}]^{T}v_{body}^{odom}
        +
        \left(
        \omega_{body}^{body}\times r_{pc}^{body}
        \right)
        \right)
        -
        \left[
        [R_{body}^{odom}]^{T}v_{body}^{odom}
        \right]_{\times}\delta R
        \\
        &=
        h(x)
        -
        \left[
        [R_{body}^{odom}]^{T}v_{body}^{odom}
        \right]_{\times}\delta R.
    \end{align}
        
    Therefore, the partial derivative of the measurement function with respect to the rotation error is
    
    \begin{equation}
    \frac{\partial h(x,\delta R)}
    {\partial \delta R}
    \bigg|_{\delta R=0}
    =
    -\left[
    [R_{body}^{odom}]^{T}v_{body}^{odom}
    \right]_{\times}.
    \end{equation} 
    
    \subsection*{Velocity}

    A small velocity error $\delta V$ changes the measurement function as

    \begin{equation}
        h(x,\delta V)
        =
        -\left(
        [R_{body}^{odom}]^{T}v_{body}^{odom}
        +
        \delta V
        +
        \left(
        \omega_{body}^{body}\times r_{pc}^{body}
        \right)
        \right).
    \end{equation}

    Using the definition of $h(x)$, this can be written as

    \begin{equation}
        h(x,\delta V)
        =
        h(x)-\delta V.
    \end{equation}

    Differentiating with respect to $\delta V$ gives

    \begin{equation}
        \frac{\partial h(x,\delta V)}
        {\partial \delta V}\bigg|_{\delta V=0}
        =
        -\mathbf{I}.
    \end{equation}

    \subsection*{Position}

    The body position does not appear in the measurement function. Therefore, adding a small position error $\delta P$ does not change the predicted measurement:

    \begin{equation}
        h(x,\delta P)=h(x).
    \end{equation}

    Differentiating with respect to $\delta P$ gives

    \begin{equation}
        \frac{\partial h(x,\delta P)}
        {\partial \delta P}\bigg|_{\delta P=0}
        =
        \mathbf{0}.
    \end{equation}

    \subsection*{Gyroscope Bias}

    The gyroscope bias is included in the bias-corrected angular velocity:

    \begin{equation}
        \omega_{body}^{body}
        =
        \omega_{meas}^{body}-b_g.
    \end{equation}

    After adding a small gyroscope-bias error $\delta b_g$, the measurement function becomes

    \begin{align}
        h(x,\delta b_g)
        &=
        -\left(
        [R_{body}^{odom}]^{T}v_{body}^{odom}
        +
        \left(
        \omega_{meas}^{body}
        -
        b_g
        -
        \delta b_g
        \right)
        \times r_{pc}^{body}
        \right)
        \\
        &=
        h(x)
        +
        \delta b_g\times r_{pc}^{body}
        \\
        &=
        h(x)
        -
        [r_{pc}^{body}]_{\times}\delta b_g.
    \end{align}

    Differentiating with respect to $\delta b_g$ gives

    \begin{equation}
        \frac{\partial h(x,\delta b_g)}
        {\partial \delta b_g}\bigg|_{\delta b_g=0}
        =
        -[r_{pc}^{body}]_{\times}.
    \end{equation}

    \subsection*{Accelerometer Bias}

    The accelerometer bias does not appear directly in the measurement function. Therefore, adding a small accelerometer-bias error $\delta b_a$ does not change the predicted measurement:

    \begin{equation}
        h(x,\delta b_a)=h(x).
    \end{equation}

    Differentiating with respect to $\delta b_a$ gives

    \begin{equation}
        \frac{\partial h(x,\delta b_a)}
        {\partial \delta b_a}\bigg|_{\delta b_a=0}
        =
        \mathbf{0}.
    \end{equation}

    \section{Numerical Jacobian} \label{appendix:NumericalJacobianCalculation}

    The analytical measurement Jacobian derived in Appendix~\ref{appendix:MeasurementJacobianDerivation} is verified using a numerical Jacobian. Each error-state variable is perturbed independently, and the corresponding change in the measurement function is calculated using the central difference approximation. This numerical Jacobian validation implementation and the following equations were developed with the help of AI assistant ChatGPT~\cite{openai_chatgpt_jacobian_2026}. The script is carefully reviewed and verified for correctness. 

    For the $i$-th error-state variable, the corresponding column of the numerical
    Jacobian is calculated as

    \begin{equation}
        H_{\mathrm{num},i}
        =
        \frac{
            h\left(x\boxplus\epsilon_i e_i\right)
            -
            h\left(x\boxplus(-\epsilon_i e_i)\right)
        }{2\epsilon_i},
        \label{eq:numericalJacobianCentralDifference}
    \end{equation}

    where, $e_i$ is the unit vector corresponding to the $i$-th error-state variable, $\epsilon_i$ is the applied perturbation, and $\boxplus$ represents the addition of the perturbation to the state. The rotation perturbation is applied on $SO(3)$ using

    \begin{equation}
        R^{+}
        =
        R\exp\left([\delta R]_{\times}\right),
    \end{equation}

    while the velocity, position, and bias perturbations are added directly to their corresponding state variables. For the selected test state and perturbation values, the analytical and numerical Jacobians were consistent. The Python implementation used for this verification is provided in Listing~\ref{lst:numericalJacobianVerification}.

    \lstinputlisting[
        language=Python,
        caption={Numerical verification of the measurement Jacobian using central differences.},
        label={lst:numericalJacobianVerification},
        frame=single,
    ]{media/config/numerical_jacobian.py}

    The analytical and numerical Jacobians obtained for the selected test state are shown in Listing~\ref{lst:numericalJacobianOutput}. Both Jacobians contain the same values, which verifies the analytical derivation for this test case.

    
    \begin{lstlisting}[
        caption={Comparison of the analytical and numerical measurement Jacobians.},
        label={lst:numericalJacobianOutput},
        basicstyle=\ttfamily\footnotesize
    ]
    
    Expected Jacobian:

    0 -0.2 -0.5 -1.0    0    0    0    0    0    0  0.8 -0.4    0    0    0
    0.2    0  1.0    0 -1.0    0    0    0    0 -0.8    0  0.5    0    0    0
    0.5 -1.0    0    0    0 -1.0    0    0    0  0.4 -0.5    0    0    0    0

    Computed numerical Jacobian:

    0 -0.2 -0.5 -1.0    0    0    0    0    0    0  0.8 -0.4    0    0    0
    0.2    0  1.0    0 -1.0    0    0    0    0 -0.8    0  0.5    0    0    0
    0.5 -1.0    0    0    0 -1.0    0    0    0  0.4 -0.5    0    0    0    0
    \end{lstlisting}

    \section{Contact Estimation Comparison} \label{sec:contactEstimationComparison}
    
    This section provides contact estimation comparison between the \gls{md} based and \gls{lf} based estimation. This is not provided in the main report because the contact estimation cannot be verified with ground truth. The ground truth can be obtained from the videos collected during the experiments. Due to the lack of time, the videos were not analyzed to get the ground truth. The comparison provided in this section uses the \gls{lf} method as a pseudo ground truth. The \gls{lf} method is used as the reference because its geometric assumption is expected to be more reliable on flat solid terrain. 
    
    The goal of this section is to show that the contact estimation problem has been worked on and evaluate the level of agreement between the \gls{md} based and \gls{lf} based contact estimates. Since the filter state is corrected using the \gls{lf} selected contacts during this comparison, the resulting \gls{md} estimates are not independent of the \gls{lf} method.
    
    For the comparison, the contact estimation is evaluated for the dataset \texttt{DATA\_04\_02}. Figure~\ref{fig:contactEstimationComparison} shows the contact estimates for each leg of the Coyote~3 rover over time. The $X$ marks the contact selected by the \gls{md} method and the $O$ marks the contact selected by the \gls{lf} method. It can be observed from the plot that, when the rover is stationary, the \gls{md} method is uncertain on which spoke is in contact. When the rover is moving, the \gls{md} method is able to select the same spoke that is selected by the \gls{lf} method. The contact estimation metrics are shown in Figure~\ref{fig:contactEstimationComparisonLegend}. The metrics are calculated over the entire dataset. Precision measures the proportion of MD-selected contacts that are also selected by the LF method. Recall measures the proportion of LF-selected contacts recovered by the MD method. The Jaccard index measures the overlap between the complete LF and MD contact sets, while the exact-match rate measures the percentage of synchronized wheel samples for which both methods select identical contact sets.
    
    \begin{figure}[!htbp]
        \centering
        \includegraphics[width=0.8\textwidth]{media/results/04_contact_estimation/Contact_Selection_Metrics_DATA_04_02.png}
        \caption[Bar plot of contact estimation metrics]{Bar plot of contact estimation metrics. }
        \label{fig:contactEstimationComparisonLegend}
    \end{figure}
    
    \begin{figure}[!htbp]
        \centering
        \includegraphics[width=\textwidth]{media/results/04_contact_estimation/Contact_Selection_LF_vs_MD_DATA_04_02.png}
        \caption[Contact estimation comparison between the \gls{md} and \gls{lf} methods.]{Contact estimation comparison between the \gls{md} and \gls{lf} methods. The contact estimates are shown for each leg of the robot over time. The \gls{lf} method is used as a pseudo ground truth for this comparison.}
        \label{fig:contactEstimationComparison}
    \end{figure}
        
    \chapter{Parameters}

    This chapter provides the OroGen and YAML configuration files of the \gls{cvoe} component. These are the configurations used in the experiments and evaluations presented in this thesis. The OroGen file defines the component's structure, including its input and output ports, operations, and properties. The YAML configuration file specifies the parameters for the component, such as filter settings, hardware configurations, and frame names. The OroGen file is provided in Listing~\ref{code:orogenfileCVOE}, while the YAML configuration file is presented in Listing~\ref{code:YAMLConfigfileCVOE}. Additionally, a sample summary of the statistical error metrics obtained from the evaluation of the \gls{cvoe} component is provided in Listing~\ref{lst:statistical-error-metrics}.
    
    \lstinputlisting[style=orogenstyle,
    caption={OroGen file of the \gls{cvoe} component},
    label={code:orogenfileCVOE}]{media/config/contact_velocity_odometry_estimator.orogen}

    \lstinputlisting[style=yamlstyle,
    caption={Contact Velocity Odometry Estimator YAML configuration file},
    label={code:YAMLConfigfileCVOE}]{media/config/contact_velocity_odometry_estimator::Task.yml}
        
    \begin{lstlisting}[style=consoleoutput, caption={Sample summary of the statistical error metrics.}, label={lst:statistical-error-metrics}]
    
    --------------------------------------------------
    Position error metrics
    --------------------------------------------------
    Max position error          : 0.92 m
    Mean position error         : 0.32 m
    Min position error          : 0.00 m
    Std dev of position error   : 0.32 m
    RMSE of position error      : 0.45 m    
    --------------------------------------------------
    Orientation Geodesic angle error metrics
    --------------------------------------------------
    Max geodesic angle error        : 9.36 degrees
    Mean geodesic angle error       : 4.80 degrees
    Min geodesic angle error        : 0.00 degrees
    Std dev of geodesic angle error : 2.65 degrees
    RMSE of geodesic angle error    : 5.49 degrees
    --------------------------------------------------
    --------------------------------------------------
    Relative Pose Error (RPE) Translation metrics
    --------------------------------------------------
    Max translation error                   : 0.05 m
    Mean translation error                  : 0.04 m
    Std dev of translation error            : 0.00 m
    RMSE of translation error               : 0.04 m

    Mean translation error per meter        : 0.15 m/m
    Std dev of translation error per meter  : 0.01 m/m
    RMSE of translation error per meter     : 0.15 m/m
    --------------------------------------------------
    Relative Pose Error (RPE) Rotation metrics
    --------------------------------------------------
    Max rotation error                  : 3.59 degrees
    Mean rotation error                 : 0.91 degrees
    Std dev of rotation error           : 0.92 degrees
    RMSE of rotation error              : 1.29 degrees

    Mean rotation error per meter       : 3.62 degrees/m
    Std dev of rotation error per meter : 3.70 degrees/m
    RMSE of rotation error per meter    : 5.18 degrees/m
    --------------------------------------------------
    \end{lstlisting}
    
    % \chapter{[in the case of using AI-assistants] Description of AI-Generated Content}
    % Describe in detail what content was generated using an AI assistant.
    % Name the AI assistant and how it was used (e.g. which prompts were used, and for which parts of the project).
    
    % ChatGPT plus was used. 
    
    % - the UML diagrams, 
    % - grammar and sentence structure corrections,
    % - latex table from excel data,
    % - latex code for equations from a rudimentary description of the equations,
    % - coherence and flow checking of the report
    % Summarizing overview of results and conclusion from the results and discussion chapter.
    
    
    
\chapter{Description of AI-Generated Content}
\label{app:aiGeneratedContent}

Artificial intelligence assistants were used as supporting tools during the implementation and writing of this thesis. The AI assistants were mainly used for language review, technical discussions, code review, debugging, equation derivation, and the preparation of evaluation scripts. All AI-generated suggestions were critically reviewed before being included in the implementation or the report. The author remains responsible for the correctness of the implementation, experiments, results, and written content.

\section{ChatGPT}

ChatGPT, developed by OpenAI, was used during different stages of the project. Its use can be grouped into the following areas.

\subsection{Language and Report Structure}

ChatGPT was used to check paragraphs for grammar, spelling, coherence, and technical clarity. It was also used to suggest improvements to sentence structure, figure captions, section outlines, and LaTeX formatting. The generated suggestions were manually reviewed and modified to maintain a consistent writing style throughout the report.

Representative prompt included: ``Check this paragraph for coherence, technical correctness, grammar, and spelling. Do not rewrite it unless restructuring is necessary.''

This assistance was used throughout the report, particularly in the Background, State of the Art, Methodology, Implementation, Evaluation and Results, and Conclusions chapters.

\subsection{Technical Discussions and Equation Verification}

ChatGPT was used to discuss and verify technical concepts related to the \gls{inekf}, rigid-body kinematics, potential contact velocity, Mahalanobis distance, contact estimation, slip rejection, and odometry evaluation metrics.

It was also used to assist with the presentation and verification of the measurement model and its Jacobian. Representative prompts included:

\begin{itemize}
    \item ``Derive the measurement Jacobian of the potential contact velocity measurement with respect to the filter error state.''
    \item ``Check the signs and coordinate frames used in the predicted and measured potential contact velocity equations.''
    \item ``Explain the rotation, velocity, position, accelerometer-bias, and gyroscope-bias blocks of the measurement Jacobian.''
\end{itemize}

\subsection{Programming Assistance}

ChatGPT was used to review and debug parts of the C++ implementation of the contact velocity odometry and its Orocos component. It was also used to assist with Python scripts for dataset processing, trajectory alignment, evaluation metric calculation, plot generation, and numerical verification. All generated or modified code was reviewed, executed, and tested using the collected datasets before being used in the project. ChatGPT was not used to generate experimental measurements or ground-truth trajectories.

Representative prompts included:

\begin{itemize}
    \item ``Review this C++ implementation of the contact velocity measurement update and identify possible frame, sign, or indexing errors.''
    \item ``Check whether the measurement covariance and innovation covariance are calculated correctly.''
    \item `` Create plantUML script for the contact velocity odometry estimator component and its input/output ports from the provided C++ code.''
\end{itemize}

\subsection{Evaluation and Result Review}

ChatGPT was used to compare written interpretations against generated plots and metric tables. This included checking whether statements about trajectory error, yaw error, translational drift, rotational drift, and contact-selection performance were supported by the corresponding results. The final interpretation and conclusions were determined by the author based on the experimental results.
Representative prompts included:

\begin{itemize}
    \item ``Compare this written interpretation with the attached trajectory plot and metric table.''
    \item ``Identify statements that are not supported by the reported evaluation metrics.''
    \item ``Check whether the dataset, figure, and table references are consistent.''
\end{itemize}

\section{Gemini CLI}

Gemini CLI, developed by Google, was used as a programming and technical reasoning assistant during the derivation and implementation of the potential contact velocity measurement Jacobian. It was used to inspect the measurement function, derive intermediate expressions, and compare the derivation with the C++ implementation.

Representative prompts included:

\begin{itemize}
    \item ``Derive the analytical Jacobian of this contact velocity measurement model with respect to the error-state variables.''
    \item ``Check whether the implemented Jacobian corresponds to the stated measurement function.''
    \item ``Identify possible sign or coordinate-frame inconsistencies in the measurement update.''
\end{itemize}

The output from Gemini CLI was treated as an intermediate suggestion. The resulting Jacobian was manually reviewed and numerically verified before being included in the implementation and Appendix~\ref{appendix:MeasurementJacobianDerivation}.

\section{Verification and Author Responsibility}

AI-generated outputs were not accepted without review. The following verification steps were applied:

\begin{itemize}
    \item Mathematical expressions were checked against the measurement model and coordinate-frame definitions.
    \item The measurement Jacobian was verified using numerical differentiation.
    \item Generated and modified code was compiled, executed, and tested on recorded datasets.
    \item Explanations of evaluation results were compared with the corresponding plots and metric tables.
    \item Literature-related statements were checked against the original cited publications.
    \item The final report text was reviewed and edited by the author.
\end{itemize}

The AI assistants were not used to collect experimental data, generate ground-truth measurements, or independently determine the final scientific conclusions. All decisions concerning the methodology, implementation, experimental design, evaluation, and conclusions were made by the author.
    
    
\end{document}
